\documentclass[10pt,a4paper]{article}

\usepackage[utf8]{inputenc}
\usepackage[T1]{fontenc}
\usepackage{geometry}
\geometry{top=0.85in,bottom=0.85in,left=0.9in,right=0.9in}
\usepackage{hyperref}
\usepackage{amsmath,amssymb}
\usepackage{mathpartir}
\usepackage{listings}
\usepackage{xcolor}
\usepackage{parskip}
\setlength{\parskip}{4pt}
\usepackage{titlesec}
\titlespacing*{\section}{0pt}{8pt}{3pt}
\titlespacing*{\subsection}{0pt}{5pt}{2pt}

\definecolor{codegray}{rgb}{0.45,0.45,0.45}
\definecolor{backcolour}{rgb}{0.96,0.96,0.94}
\lstdefinestyle{mystyle}{
    backgroundcolor=\color{backcolour},
    basicstyle=\ttfamily\scriptsize,
    breaklines=true, keepspaces=true,
    frame=single, framerule=0.3pt,
    rulecolor=\color{codegray!40},
    numbers=none, showstringspaces=false, tabsize=2
}
\lstset{style=mystyle}

\title{\textbf{Rusty 2 Interpreter --- Phase 2 Report}}
\author{Artur Krystopchuk \\ \small ist1104145}
\date{\today}

\begin{document}
\maketitle\vspace{-10pt}

% ------------------------------------------------------------------
\section{Architecture}

The implementation uses the \emph{Visitor} pattern: \texttt{TypeCheckVisitor} and
\texttt{EvalVisitor} each traverse the AST, threading a scoped \texttt{Environment<T>}.
Three private helpers in \texttt{TypeCheckVisitor} carry most of the logic:
\texttt{unfold} resolves type-alias names, \texttt{isSubtype} implements $A <: B$, and
\texttt{leastUpperBound} computes the join of two types.

% ------------------------------------------------------------------
\section{Static Type Checker}

\paragraph{Subtyping.}
\texttt{isSubtype} covers reflexivity, then three structural cases.
\emph{Arrows} are contravariant in the domain and covariant in the codomain.
\emph{References} (\texttt{ref} and \texttt{\&}) are \emph{invariant}: both
$A <: B$ and $B <: A$ must hold, preventing variance through mutable
locations.
\emph{Enums} combine width and depth subtyping in a single check: for $E_1 <: E_2$,
every tag in $E_1$ must appear in $E_2$, and for each shared tag the field tuple of
$E_1$ must be at least as long as $E_2$'s, with each shared prefix field a subtype:
$$\frac{(\text{all } i \in 1\dots k) \ A_i = (C_1, \dots, C_{a_i}, \dots, C_{b_i}) <: B_i = (D_1, \dots, D_{a_i}) \quad C_p <: D_p \ (\text{all } p \in p\dots a_i)}{\text{enum } \{l_1 : A_1, \dots, l_k : A_k\} <: \text{enum } \{l_1 : B_1, \dots, l_k : B_k, \dots, l_n : B_n\}}$$


\paragraph{Least upper bound.}
\texttt{leastUpperBound} is called in four situations:
\texttt{visitIf} and \texttt{visitMatch} use it to unify the types of
the two (or more) branches into a single return type.
\texttt{visitEq} and \texttt{visitDif} use it purely to \emph{check}
that the two operands are type-compatible before returning \texttt{bool}.
For enums the LUB merges tag sets and trims each shared tag's field list to the
longest common equal-type prefix.

\paragraph{Pattern matching.}
\texttt{visitMatch} checks that each pattern tag belongs to the matched enum, that
bound-variable counts do not exceed field counts, and that the match is exhaustive
(every tag in the enum type must appear in some branch).

% ------------------------------------------------------------------
\section{Recursive Types (Challenge 1)}

\texttt{visitTypeDefinition} simply inserts name-to-type bindings into a new scope.
\texttt{ASTTypeId} represents a not-yet-resolved alias; \texttt{unfold} resolves it
by looking up the name and recursing, so a self-referential definition such as
\texttt{type Btree = enum \{ \#nil:(unit), \#node:(Btree,int,Btree) \}} works because
the inner \texttt{Btree} is only expanded when a structural comparison actually needs it.

Recursive \emph{functions} require an explicit type annotation.
In \texttt{visitLet}, if an annotation is present the declared type is placed in
the environment \emph{before} the RHS is evaluated, letting the body refer back to
the same identifier during type-checking.
Without an annotation the binding is inserted only after the RHS, so unannotated
recursive functions are correctly rejected.

% ------------------------------------------------------------------
\section{Stack-Allocated Mutable Variables (Challenge 2)}

\paragraph{Type-checking.}
\texttt{ASTTypeMut(A)} is an internal wrapper never visible in the concrete syntax.
A \texttt{let mut x = e} binding wraps the inferred or declared type as
$\mathit{mut}(A)$ before storing it in the type environment.
\texttt{visitId} implicitly unwraps $\mathit{mut}(A)$ on reads, giving automatic
R-value dereferencing.
Assignments are routed through \texttt{typeOfAssignable}, which accepts
\texttt{ASTTypeMut}, \texttt{ASTTypeRef}, and \texttt{ASTTypeAddr} as valid targets.

\paragraph{Evaluation.}
Each \texttt{let mut} binding calls \texttt{memory.push(v)}, receiving a
\texttt{VirtualAddress} stored as a \texttt{VMutAddress} in the value environment.
\texttt{visitId} silently reads through any \texttt{VMutAddress} via
\texttt{memory.read}.
The \texttt{finally} block of \texttt{visitLet} calls \texttt{memory.pop()} once
per mutable binding, enforcing the LIFO (Last-In, First-Out).

% ------------------------------------------------------------------
\section{Explicit Aliases (Challenge 3)}

\texttt{ASTTypeAddr(A)} is the programmer-visible alias type \texttt{\&(A)}.
\texttt{visitAddr} in the type checker verifies that the target identifier is bound
to \texttt{ASTTypeMut(A)} and returns \texttt{ASTTypeAddr(A)}.
Scope-escape is prevented statically: a function whose return type is \texttt{\&(A)}
but that takes the alias of an inner \texttt{let mut} variable will fail because
the \texttt{ASTTypeMut} binding no longer exists in the type environment at the
call site.
Invariant subtyping on \texttt{ASTTypeAddr} closes any variance loophole.

At runtime, \texttt{visitAddr} retrieves the raw \texttt{VirtualAddress} from the
\texttt{VMutAddress} in the environment.
Dereference and assignment handle both \texttt{VCell} (heap) and
\texttt{VirtualAddress} (stack alias) uniformly, so the same \texttt{!} and
\texttt{:=} operators work across all reference kinds.

\end{document}
